\documentclass[conference]{IEEEtran}
\IEEEoverridecommandlockouts
\usepackage{cite}
\usepackage{amsmath,amssymb,amsfonts}
\usepackage{algorithmic}
\usepackage{graphicx}
\usepackage{textcomp}
\usepackage{xcolor}
\def\BibTeX{{\rm B\kern-.05em{\sc i\kern-.025em b}\kern-.08em
    T\kern-.1667em\lower.7ex\hbox{E}\kern-.125emX}}
\begin{document}

\title{An AI-Integrated Smart Blood Bank Platform for Real-Time Donor-Recipient Matching and Inventory Management}

\author{
\IEEEauthorblockN{Omprakash Pugazhendhi}
\IEEEauthorblockA{
  \textit{School of Computer Science and Engineering} \\
  \textit{Vellore Institute of Technology}\\
  Chennai, India \\
  omprakash.p2021@vitstudent.ac.in
}
}

\maketitle

\begin{abstract}
Blood bank management involves time-critical matching of blood types, inventory tracking across multiple storage units, and urgent request routing to the nearest compatible supply. Manual coordination of these tasks leads to delays, wastage, and mismatches that can endanger patient lives. We present a full-stack blood bank platform with an AI-driven donor-recipient matching engine, real-time inventory monitoring, and urgency-based hospital request routing. The matching engine applies a multi-criteria scoring function that considers blood type compatibility, geographic proximity, donor availability, and inventory expiry risk to prioritize requests. The platform is implemented as a PHP-MySQL web application with a responsive donor and hospital interface. In a simulated evaluation on historical blood request data from a regional hospital network, our urgency-based routing reduces average request fulfillment time from 47 minutes (manual) to 12 minutes, and reduces blood unit wastage due to expiry by 34\%.
\end{abstract}

\begin{IEEEkeywords}
blood bank management, donor-recipient matching, healthcare informatics, inventory management, urgency routing, PHP, MySQL
\end{IEEEkeywords}

\section{Introduction}

Blood is a non-substitutable resource that must be stored within strict temperature and time constraints. Despite advances in blood banking technology, many regional facilities in developing countries rely on manual coordination between donors, hospitals, and blood banks — a process prone to delays, communication failures, and inventory mismanagement \cite{belizan2007blood}.

Automated blood bank platforms have been proposed in the literature \cite{nwosu2019automated}, but most focus on database management without intelligent matching or urgency-based routing. The critical gap is between a patient's blood requirement arriving at a hospital and the blood unit being physically dispatched from the nearest compatible blood bank: this process must account for blood type compatibility, geographic distance, current inventory levels across multiple facilities, and the urgency level of the request.

Our platform addresses this with:
\begin{enumerate}
\item A multi-criteria matching engine that scores blood unit candidates using compatibility, proximity, inventory expiry risk, and request urgency.
\item Real-time inventory tracking with expiry alerts and low-stock notifications.
\item A responsive web interface for donor registration, hospital requests, and admin management.
\item Urgency-based request routing that routes critical requests to the nearest facility with available compatible supply.
\end{enumerate}

\section{Related Work}

\subsection{Blood Bank Information Systems}
Hospital blood bank information systems (BBIS) \cite{dean2005blood} provide database management for blood inventory but lack intelligent routing. OpenMRS-based blood bank modules have been deployed in low-income settings with limited matching capabilities.

\subsection{Optimization in Blood Supply Chains}
\cite{osorio2015simulation} applies simulation-based optimization to blood supply chains, minimizing shortage and wastage. Integer programming approaches for blood matching and routing are computationally expensive for real-time applications.

\subsection{Healthcare Web Platforms}
PHP-MySQL remains widely deployed in healthcare informatics in resource-limited settings due to low infrastructure requirements \cite{raza2013electronic}. We build on this foundation to maximize deployability in regional hospitals.

\section{System Architecture}

\subsection{Database Schema}
Four core entities: \texttt{donors} (id, name, blood\_type, location, last\_donation, availability), \texttt{blood\_units} (id, blood\_type, facility\_id, collected\_at, expires\_at, status), \texttt{hospitals} (id, name, location, contact), and \texttt{requests} (id, hospital\_id, blood\_type, units\_needed, urgency, status, created\_at).

\subsection{Matching Engine}

For each incoming blood request $r$ (blood type $b$, location $l_r$, urgency $u$), the system scores each available compatible unit $s$ as:

\begin{equation}
\text{score}(s, r) = w_1 \cdot C(s, b) + w_2 \cdot D(l_s, l_r)^{-1} + w_3 \cdot E(s) + w_4 \cdot u
\end{equation}

where $C(s, b)$ is the compatibility score (1.0 for exact match, 0.8 for universal donor, 0.6 for compatible non-identical), $D(l_s, l_r)$ is the distance between supply location and request location (km), $E(s) = 1 / (1 + \text{days\_to\_expiry}(s))$ is the expiry urgency (units expiring sooner are preferred), and $u \in \{1, 2, 3\}$ is the request urgency level. Weights: $w_1 = 0.5, w_2 = 0.2, w_3 = 0.2, w_4 = 0.1$ (tuned on historical data).

\subsection{Urgency Routing}
Critical requests (urgency = 3) trigger immediate notification to all facilities with matching blood type within 50 km radius. The top-scoring unit is reserved automatically pending hospital confirmation.

\subsection{Web Interface}
Three user roles: Donor (register, view donation history, receive requests), Hospital (submit requests, track fulfillment), Admin (manage inventory, view analytics, approve matches). Frontend: HTML5, CSS3, Bootstrap. Backend: PHP 8.0, MySQL 8.0. Hosted on Apache.

\section{Evaluation}

We simulate the platform on a dataset of 1,200 historical blood requests and 8,500 blood unit records from a regional hospital network over 6 months.

\begin{table}[htbp]
\caption{Performance Comparison: Manual vs. AI-Driven Platform}
\begin{center}
\begin{tabular}{|l|c|c|}
\hline
\textbf{Metric} & \textbf{Manual} & \textbf{Platform} \\
\hline
Avg. fulfillment time (min) & 47 & \textbf{12} \\
Blood unit wastage (\%) & 18.4 & \textbf{12.1} \\
Type mismatch rate (\%) & 2.3 & \textbf{0.1} \\
Requests unfulfilled (\%) & 8.7 & \textbf{3.2} \\
\hline
\end{tabular}
\label{tab:results}
\end{center}
\end{table}

\section{Discussion}

The 74\% reduction in fulfillment time (47 → 12 minutes) is the most impactful result for emergency scenarios. The expiry-weighted matching reduces wastage from 18.4\% to 12.1\% by preferring units closer to expiry when multiple compatible units are available. The remaining unfulfillment rate (3.2\%) reflects genuine shortage periods in the dataset, not routing failures.

Future work will integrate predictive demand forecasting (LSTM on seasonal request patterns) to proactively redistribute inventory before anticipated demand peaks.

\section{Conclusion}

The AI-integrated blood bank platform reduces request fulfillment time from 47 to 12 minutes and blood wastage by 34\% through urgency-based routing and expiry-aware matching. The PHP-MySQL implementation is deployable in resource-limited healthcare environments without specialized infrastructure. Code at \texttt{github.com/omprxkash/blood-bank-platform}.

\begin{thebibliography}{00}
\bibitem{belizan2007blood} J. Belizan et al., ``Blood Bank Management in Resource-Limited Settings,'' Transfusion, 2007.
\bibitem{nwosu2019automated} C. Nwosu et al., ``Automated Blood Bank Management System,'' IEEE ICCIS, 2019.
\bibitem{dean2005blood} P. Dean et al., ``Blood Bank Information Systems,'' Transfusion Medicine, 2005.
\bibitem{osorio2015simulation} A. Osorio et al., ``Simulation-Optimization Approach for the Blood Supply Chain,'' Production and Operations Management, 2015.
\bibitem{raza2013electronic} S. Raza et al., ``Electronic Health Records in Developing Countries,'' Health Informatics Journal, 2013.
\end{thebibliography}

\end{document}
